\فصل{مبانی نظری}

در این فصل، مبانی نظری و مفاهیم پایه‌ای مورد استفاده در این پژوهش ارائه می‌شود.
ابتدا مدل ریاضی رسمی مسئله زمان‌بندی کارگاهی معرفی می‌شود.
سپس مفاهیم بهینه‌سازی چندهدفه شامل غلبه پارتو، شاخص‌های ارزیابی
و الگوریتم \lr{NSGA-II} مرور می‌شوند.
در نهایت، مدل زیرساخت رایانشی لبه-مه-ابر
شامل ویژگی‌های گره‌ها، مدل ارتباطی و توابع هدف، تشریح می‌شود.

\قسمت{مدل ریاضی مسئله زمان‌بندی کارگاهی}

\زیرقسمت{تعریف رسمی}

مسئله زمان‌بندی کارگاهی
(\lr{JSSP})
با $n$ کار و $m$ ماشین تعریف می‌شود؛
$n$ تعداد کارها و $m$ تعداد ماشین‌ها است.
مسئله با پارامترهای زیر تعریف می‌شود~\cite{pinedo2016scheduling}:

\شروع{فقرات}
\فقره مجموعه کارها: $J = \set{J_1, J_2, \dots, J_n}$
\فقره مجموعه ماشین‌ها: $M = \set{M_1, M_2, \dots, M_m}$
\فقره هر کار $J_j$ شامل $m$ عملیات مرتب $O_{j,1}, O_{j,2}, \dots, O_{j,m}$ است
\فقره هر عملیات $O_{j,k}$ باید روی ماشین $\mu(j,k)$ با زمان پردازش $p_{j,k}$ اجرا شود
\پایان{فقرات}

\شروع{تعریف}[مسئله زمان‌بندی کارگاهی]
با داشتن مجموعه‌های $J$ و $M$ و ماتریس زمان‌های پردازش $P = [p_{j,k}]$
و تابع نگاشت ماشین $\mu: J \times \set{1,\dots,m} \to M$،
هدف یافتن برنامه زمان‌بندی مجاز $\sigma$ است
که زمان تکمیل کل $C_{\max}(\sigma) = \max_{j} C_j$ را کمینه کند،
به‌طوری‌که محدودیت‌های زیر برآورده شوند:

\شروع{شمارش}
\فقره \مهم{محدودیت تقدم}: عملیات $O_{j,k}$ باید پیش از $O_{j,k+1}$ تکمیل شود.
\فقره \مهم{محدودیت ماشین}: هر ماشین در هر لحظه حداکثر یک عملیات را اجرا می‌کند.
\فقره \مهم{عدم وقفه}: هر عملیات پس از شروع بدون وقفه تا پایان اجرا می‌شود.
\پایان{شمارش}
\پایان{تعریف}


\زیرقسمت{نمادگذاری}

در جدول~\رجوع{جدول:نمادها} نمادهای اصلی مورد استفاده در این پایان‌نامه فهرست شده‌اند.

\شروع{لوح}[ht]
\تنظیم‌ازوسط
\شرح{نمادهای اصلی مورد استفاده}

\شروع{جدول}{|c|r|}
\خط‌پر
\سیاه نماد & \سیاه توضیح \\
\خط‌پر \خط‌پر
$n$ & تعداد کارها \\
$m$ & تعداد ماشین‌ها \\
$J_j$ & کار $j$ام \\
$M_i$ & ماشین $i$ام \\
$O_{j,k}$ & عملیات $k$ام از کار $j$ \\
$p_{j,k}$ & زمان پردازش عملیات $O_{j,k}$ \\
$C_j$ & زمان تکمیل کار $j$ \\
$C_{\max}$ & زمان تکمیل کل (\lr{Makespan}) \\
$E_{\text{total}}$ & مصرف انرژی کل \\
$R$ & قابلیت اطمینان سامانه \\
\خط‌پر
\پایان{جدول}

\برچسب{جدول:نمادها}
\پایان{لوح}


\قسمت{بهینه‌سازی چندهدفه}

\زیرقسمت{تعریف مسئله چندهدفه}

در بهینه‌سازی چندهدفه\پاورقی{Multi-Objective Optimization}
هدف بهینه‌سازی هم‌زمان چند تابع هدف متناقض است~\cite{deb2001multi}.
یک مسئله بهینه‌سازی چندهدفه به‌صورت زیر تعریف می‌شود:

\begin{equation}
\min_{\mathbf{x} \in \Omega} \mathbf{f}(\mathbf{x}) = (f_1(\mathbf{x}), f_2(\mathbf{x}), \dots, f_k(\mathbf{x}))
\end{equation}

که $\Omega$ فضای مجاز و $k$ تعداد توابع هدف است.

\شروع{تعریف}[غلبه پارتو]
راه‌حل $\mathbf{x}$ بر راه‌حل $\mathbf{y}$ غلبه دارد ($\mathbf{x} \prec \mathbf{y}$)
اگر و تنها اگر برای همه $i$: $f_i(\mathbf{x}) \le f_i(\mathbf{y})$
و برای حداقل یک $i$: $f_i(\mathbf{x}) < f_i(\mathbf{y})$.
\پایان{تعریف}

\شروع{تعریف}[جبهه پارتو]
مجموعه تمام راه‌حل‌های غیرمغلوب، \مهم{جبهه پارتو}\پاورقی{Pareto front} نامیده می‌شود.
\پایان{تعریف}


\زیرقسمت{شاخص‌های کیفیت}

برای ارزیابی کیفیت مجموعه راه‌حل‌های پارتو از شاخص‌های زیر استفاده می‌شود:

\شروع{فقرات}
\فقره
\مهم{ابرحجم}\پاورقی{Hypervolume (HV)}:
حجم فضایی که توسط مجموعه راه‌حل‌های غیرمغلوب و یک نقطه مرجع محصور می‌شود~\cite{zitzler2003performance}.
مقدار بیشتر ابرحجم نشان‌دهنده کیفیت بهتر است.
در این پایان‌نامه، نقطه مرجع مشترکی برای هر نمونه مسئله از اجتماع تمام نقاط تولیدشده توسط همه الگوریتم‌ها ساخته می‌شود.
\فقره
\مهم{فاصله نسلی معکوس}\پاورقی{Inverted Generational Distance (IGD)}:
میانگین فاصله هر نقطه از جبهه مرجع تا نزدیک‌ترین نقطه در مجموعه راه‌حل~\cite{zitzler2003performance}.
مقدار کمتر \lr{IGD} نشان‌دهنده کیفیت بهتر است.
در این پایان‌نامه، جبهه مرجع تجربی از اجتماع جبهه‌های غیرمغلوب تمام الگوریتم‌ها ساخته می‌شود.
\پایان{فقرات}


\زیرقسمت{الگوریتم \lr{NSGA-II}}

الگوریتم \lr{NSGA-II}\پاورقی{Non-dominated Sorting Genetic Algorithm II}
یکی از پرکاربردترین الگوریتم‌های بهینه‌سازی تکاملی چندهدفه است
که توسط دب و همکاران ارائه شده است~\cite{deb2002fast}.
این الگوریتم دارای ویژگی‌های زیر است:

\شروع{شمارش}
\فقره \مهم{مرتب‌سازی غیرمغلوب سریع}: جمعیت به لایه‌های غیرمغلوب دسته‌بندی می‌شود.
بر راه‌حل‌های لایه اول، هیچ راه‌حل دیگری غلبه نمی‌کند.
\فقره \مهم{فاصله ازدحام}\پاورقی{Crowding Distance}: معیاری برای حفظ تنوع در هر لایه.
راه‌حل‌هایی که در نواحی کم‌جمعیت‌تر قرار دارند، ارجحیت بیشتری دارند.
\فقره \مهم{انتخاب تورنمنتی}: والدین بر اساس ترکیب رتبه غلبه و فاصله ازدحام انتخاب می‌شوند.
\فقره \مهم{حفظ نخبگان}: ادغام جمعیت والدین و فرزندان و انتخاب بهترین‌ها.
\پایان{شمارش}


\قسمت{مدل زیرساخت رایانشی لبه-مه-ابر}

\زیرقسمت{معماری سه‌لایه}

زیرساخت رایانشی مورد استفاده در این پایان‌نامه شامل سه‌لایه است:

\شروع{شمارش}
\فقره \مهم{لایه لبه}\پاورقی{Edge}: شامل ۴ تا ۶ گره با توان پردازشی \lr{120 MIPS}،
توان مصرفی فعال ۱۸ وات، توان بی‌کاری ۴.۵ وات، و نرخ خرابی $\lambda = 0.020$.
\فقره \مهم{لایه مه}\پاورقی{Fog}: شامل ۳ تا ۴ گره با توان پردازشی \lr{320--340 MIPS}،
توان مصرفی فعال ۴۲ وات، توان بی‌کاری ۹ وات، و نرخ خرابی $\lambda = 0.010$.
\فقره \مهم{لایه ابر}\پاورقی{Cloud}: شامل ۲ گره با توان پردازشی \lr{760--800 MIPS}،
توان مصرفی فعال ۹۵ وات، توان بی‌کاری ۲۲ وات، و نرخ خرابی $\lambda = 0.005$.
\پایان{شمارش}

شکل~\رجوع{شکل:توپولوژی} توپولوژی شبکه زیرساخت را نشان می‌دهد.

\شروع{شکل}[ht]
\centerimg{topology.pdf}{12cm}
\شرح{توپولوژی شبکه لبه-مه-ابر و پیوندهای ارتباطی میان لایه‌ها.}
\برچسب{شکل:توپولوژی}
\پایان{شکل}

همان‌طور که در شکل~\رجوع{شکل:توپولوژی} دیده می‌شود، ساختار شبکه به‌صورت سلسله‌مراتبی
از لایه لبه به مه و سپس ابر سازمان‌دهی شده است. لایه لبه نزدیک‌ترین نقطه به منبع تولید داده است
و برای تصمیم‌های سریع مناسب است؛ لایه مه نقش میانجی را برای تجمیع بار، کاهش ازدحام لینک‌های بالادست
و مدیریت نوسان بار ایفا می‌کند؛ و لایه ابر ظرفیت پردازشی بالاتر برای پردازش‌های سنگین‌تر را فراهم می‌سازد.


\زیرقسمت{مدل ارتباطات}

ارتباط بین گره‌ها از طریق پیوندهایی با پهنای باند و تأخیر انتشار مشخص انجام می‌شود.
شبیه‌سازی با استفاده از \lr{YAFS}\پاورقی{Yet Another Fog Simulator}~\cite{lera2019yafs} صورت گرفته است.
برای هر انتقال داده روی پیوند $e$:

\begin{equation}
t_{\text{tx}}(e) = \frac{\text{bytes}}{BW_e \times 10^6}
\end{equation}

\begin{equation}
\text{lat}_e = t_{\text{tx}}(e) + PR_e
\end{equation}

که $BW_e$ پهنای باند (بر حسب مگابیت بر ثانیه) و $PR_e$ تأخیر انتشار است.

در این پژوهش، ارتباط بین سه لایه فقط انتقال داده خام نیست، بلکه بخشی از هزینه تصمیم‌گیری زمان‌بندی است.
هر جابه‌جایی عملیات بین گره‌ها می‌تواند به افزایش تأخیر انتهابه‌انتها و انرژی ارتباطی منجر شود؛
ازاین‌رو در ارزیابی هر کروموزوم، مسیرهای ارتباطی متناظر با نگاشت عملیات نیز هم‌زمان محاسبه می‌شوند.

\زیرقسمت{منطق تخصیص عملیات بین لایه‌ها}

در روش پیشنهادی، نگاشت عملیات به لایه‌ها توسط بخش \lr{Machine Map} کروموزوم تعیین می‌شود
و بهترین نگاشت از طریق بهینه‌سازی هم‌زمان سه هدف انتخاب می‌گردد. بااین‌حال، منطق تصمیم‌گیری
به‌صورت مفهومی به شکل زیر قابل تفسیر است:

\شروع{فقرات}
\فقره \مهم{ارسال به لایه لبه}: زمانی که عملیات به تأخیر کم حساس است، حجم داده کوچک‌تر است
و اجرای محلی می‌تواند از هزینه رفت‌وبرگشت شبکه جلوگیری کند.
\فقره \مهم{ارسال به لایه مه}: زمانی که نیاز به موازنه بین زمان پاسخ و مصرف انرژی وجود دارد
یا لایه لبه با محدودیت ظرفیت مواجه است؛ لایه مه نقطه میانی برای تجمیع و توزیع بار محسوب می‌شود.
\فقره \مهم{ارسال به لایه ابر}: زمانی که عملیات محاسباتی سنگین‌تر است و بهره‌گیری از توان پردازشی بالا
اثر غالب دارد، حتی اگر هزینه ارتباطی و مصرف انرژی افزایش یابد.
\پایان{فقرات}

بنابراین انتخاب نهایی مقصد اجرا یک قاعده ثابت از پیش‌تعریف‌شده نیست، بلکه حاصل موازنه پویا
میان سه معیار زمان، انرژی و قابلیت‌اطمینان در چارچوب بهینه‌سازی چندهدفه است.


\زیرقسمت{توابع هدف}

در این پایان‌نامه، مسئله در قالب کمینه‌سازی هم‌زمان سه تابع هدف صورت‌بندی می‌شود:

\مهم{۱. زمان تکمیل کل} ($C_{\max}$):
زمان اتمام آخرین عملیات.
اگر عملیات $O_{j,k}$ روی گره $n$ اجرا شود، زمان پردازش آن برابر $t_{\text{proc}} = p_{j,k} / IPT_n$ است که در آن $IPT_n$ توان پردازشی گره $n$ است.

\مهم{۲. مصرف انرژی کل} ($E_{\text{total}}$):
\begin{equation}
E_{\text{total}} = E_{\text{comp}} + E_{\text{idle}} + E_{\text{comm}}
\end{equation}
که:
\شروع{فقرات}
\فقره $E_{\text{comp}} = \sum_{\text{op}} P_{\text{active}}(\text{node}(\text{op})) \times t_{\text{proc}}(\text{op})$:
انرژی پردازش فعال
\فقره $E_{\text{idle}} = \sum_{n} P_{\text{idle}}(n) \times \max(0, C_{\max} - \text{busy}_n)$:
انرژی بی‌کاری گره‌ها
\فقره $E_{\text{comm}} = \sum_{\text{link}} P_{\text{tx}}(\text{link}) \times t_{\text{tx}}(\text{link})$:
انرژی ارتباطات
\پایان{فقرات}

\مهم{۳. عدم‌قابلیت‌اطمینان} ($1 - R$):
\begin{equation}
R = \prod_{\text{op}} \left[ R_{\text{node}}(\text{op}) \times \prod_{e \in \text{path}(\text{op})} R_{\text{link}}(e) \right]
\end{equation}
که:
\شروع{فقرات}
\فقره $R_{\text{node}}(\text{op}) = \exp(-\lambda_{\text{node}} \times t_{\text{proc}}(\text{op}))$
\فقره $R_{\text{link}}(e) = \exp(-\lambda_{\text{link}} \times \text{lat}_e)$
\پایان{فقرات}

بنابراین، بردار هدف برای کمینه‌سازی عبارت است از:
$\mathbf{f} = (C_{\max}, \; E_{\text{total}}, \; 1 - R)$


\قسمت{جمع‌بندی فصل}

در این فصل، مبانی نظری لازم برای ادامه پژوهش ارائه شد.
ابتدا مدل رسمی مسئله زمان‌بندی کارگاهی و نمادگذاری آن تعریف گردید،
سپس مفاهیم کلیدی بهینه‌سازی چندهدفه و الگوریتم \lr{NSGA-II} تشریح شد
و در پایان، مدل زیرساخت رایانشی لبه-مه-ابر و توابع هدف سه‌گانه پژوهش بیان شد.
فصل بعدی با تکیه بر این مبانی، به مرور نظام‌مند کارهای پیشین می‌پردازد.
